% !TEX root = ../main.tex
\begin{abstract}
Long-term timber harvest schedules produced by linear programming (LP) are routinely presented to decision-makers as credible statements about the future, and their shadow prices are used to value constraints, land, and policy trade-offs. In an unpublished 1991 PhD thesis, P. J. Daugherty showed that LP forest-planning models solved as \emph{open-loop} formulations admit \emph{dynamically inconsistent} plans---plans that fail Bellman's principle of optimality, because a future planner re-solving from the realized forest state under the same objectives would not follow the plan's tail. Such plans are not a credible basis for policy, and trade-off analysis derived from them is biased. The thesis was never peer-reviewed and is effectively inaccessible (archival hard copy only, \textasciitilde\$200 custom print-on-demand), so the phenomenon it documents has been repeatedly rediscovered. This paper is a substantive replication of previously reported results, not new science. We reproduce Daugherty's Umpqua National Forest FORPLAN case study in the open-source \texttt{ws3} wood-supply model (Paradis 2026) from the real Umpqua planning documents, and reproduce all of the thesis's headline findings on a sequential-replanning simulation: dynamic inconsistency occurs in every simulated condition (100\% of the experiment-grid cells); it is invariant to the discount rate across 0--6\%; tighter harvest-flow constraints increase it; disequilibrium forest structure with negatively valued strata drives it; and on the all-mature landbase the open-loop even-flow plan delivers roughly half its promised volume under replanning. Our only novel content is the open, reproducible software stack and a citable reference for a result the field should stop rediscovering.
\end{abstract}
